(1) 将电子速度复幅记为 \(\boldsymbol v_\omega\)，方程给 \(m(\nu-i\omega)\boldsymbol v_\omega=-e\boldsymbol E_\omega\)。电子电流密度为 \(\boldsymbol J=-ne\boldsymbol v_\omega\)，因此 \(\boldsymbol J_\omega=\sigma(\omega)\boldsymbol E_\omega\)，\(\sigma(\omega)=ne^2/[m(\nu-i\omega)]\)。符号自洽地表明当 \(\nu>0\) 时实电导正，电磁场平均向电子作正功。
(2) 对固定离子背景，定义 \(\boldsymbol D=\varepsilon_0\boldsymbol E+\boldsymbol P\)，电子导电响应等价于 \(\boldsymbol J=\partial_t\boldsymbol P=-i\omega\boldsymbol P\)。所以 \(\varepsilon(\omega)=\varepsilon_0+i\sigma(\omega)/\omega=\varepsilon_0[1-\omega_p^2/\{\omega(\omega+i\nu)\}]\)，其中 \(\omega_p^2=ne^2/(\varepsilon_0m)\)。在无碰撞极限 \(\varepsilon_r=1-\omega_p^2/\omega^2\)，纵向电荷振荡在 \(\varepsilon_r=0\) 时发生。
(3) Maxwell 横波方程给 \(k^2=\mu_0\varepsilon(\omega)\omega^2\)，且 \(c^{-2}=\mu_0\varepsilon_0\)。令 \(\nu=0\)，得到 \(\omega^2=\omega_p^2+c^2k^2\)，只有 \(\omega>\omega_p\) 才有实数传播波数。\(\omega=\omega_p\) 时 \(k=0\)，是理想模型的截止边界，并非高传播速度的自由波。
(4) 对 \(\omega>\omega_p\)，相速度 \(v_{\rm ph}=\omega/k=c/\sqrt{1-\omega_p^2/\omega^2}>c\)，群速度 \(v_g=d\omega/dk=c\sqrt{1-\omega_p^2/\omega^2}<c\)，且 \(v_{\rm ph}v_g=c^2\)。相速度超过 \(c\) 不表示信息超光速；若 \(0<\omega<\omega_p\)，\(k=i\sqrt{\omega_p^2-\omega^2}/c\)，场变为倏逝衰减而非携能进入无限均匀等离子体的传播波。
